\documentclass[a4paper,12pt]{article}
\usepackage[inner=20truemm,outer=20truemm]{geometry}
\usepackage[T1]{fontenc}
\usepackage[utf8]{inputenc}
\usepackage[magyar]{babel}
\usepackage{palatino}
\usepackage{amsmath}   
\usepackage{amssymb}
\setlength{\parindent}{0pt}
% struktogram rajzolásához be kell tölteni a .tex forrásfájllal azonos könyvtárba helyezett stuki.sty-t
\usepackage{stuki}
% specifikációban címkézett változó írásához, A=(x:Z,l:L) alakban
\newcommand{\lv}{{:}}

\begin{document}
	\author{Zsolt Borsi}	
\bigskip
\bigskip

\section{Programkonstrukciók struktogramjai}
\subsection{Szekvencia}
\begin{stuki*}[3cm]{S}
  \stm{S_1}
  \stm{S_2}
\end{stuki*}
\smallskip

\subsection{Elágazás}
\begin{stuki*}[3cm]{IF}
\begin{CASE}{1}{3}
    \WHEN{\stm{\pi_1}}
      \stm{S_1}
    \WHEN*{\stm{\cdots}}
      \stm{\cdots}
    \WHEN{\stm{\pi_n}}
      \stm{S_n}      
\end{CASE}
\end{stuki*}
\bigskip

\begin{stuki}[3cm]
\begin{IF}{1}{\stm{\pi}}
   \stm{S_1}
   \ELSE
     \stm{S_2}
\end{IF}
\end{stuki}
\smallskip

\subsection{Ciklus}
\begin{stuki*}[3cm]{DO}
\begin{WHILE}{1}{\stm{\pi}}
 \stm{S_0}
\end{WHILE}
\end{stuki*}
\newpage

%%%%%%%%%%%%%%%%%%%%%%%%%%%%%%%%%%%%%%%%%%%%%%%%%%%%%%%%%%%%%%%%%%%%%%%%%%%%%%%%%%%%%%%%%%%%%%%%%%%%%%%%%%%%%%%

\section{Néhány intervallumos programozási tétel}

%%% ÖSSZEGZÉS TÉTEL
\subsection{Összegzés}
Legyen adott az $f\colon \mathbb{Z}\rightarrow\mathbb{Z}$ függvény. Egy adott $[m..n]\subset \mathbb{Z}$ intervallumban összegezzük az $f$ függvény értékeit.\\

%%% specifikáció
$A = (m\lv\mathbb{Z}, n\lv\mathbb{Z}, s\lv\mathbb{Z})$\\[2pt]
$B = (m'\lv\mathbb{Z}, n'\lv\mathbb{Z})$\\[2pt]
$Q = (m = m' \land n = n' \land m \leqslant n+1)$\\[2pt]
$R = (Q \land s = \sum\limits_{i=m}^{n}f(i))$
%%% struktogram
\medskip
\begin{stuki}
  \stm{k,s := m-1,0}
  \begin{WHILE}{2}{\stm{k \neq n}}
      \stm{s:= s + f(k+1)}
      \stm{k:= k+1}
  \end{WHILE}
\end{stuki}

%%% SZÁMLÁLÁS TÉTEL
\subsection{Számlálás}
Legyen $\beta\colon \mathbb{Z}\rightarrow\mathbb{L}$ egy adott függvény. Számoljuk meg hány helyen igaz $\beta$ egy adott $[m..n]\subset \mathbb{Z}$ intervallumban.\\

%%% specifikáció
$A = (m\lv\mathbb{Z}, n\lv\mathbb{Z},d\lv\mathbb{Z})$\\[2pt]
$B = (m'\lv\mathbb{Z}, n'\lv\mathbb{Z})$\\[2pt]
$Q = (m = m' \land n = n' \land m \leqslant n+1)$\\[2pt]
$R = (Q \land d = \sum\limits_{i=m}^{n}\chi(\beta(i)))$\\[2pt]
ahol $\chi\colon \mathbb{L}\rightarrow \{0,1\}$ úgy hogy $\chi(igaz) = 1$ és $\chi(hamis) = 0$.
%%% struktogram
\medskip
\begin{stuki}
  \stm{k,d := m-1,0}
  \begin{WHILE}{3}{\stm{k \neq n}}
    \begin{IF}{1}{\stm{\beta(k+1)}}
      \stm{d:= d+1}
      \ELSE
      \stm{SKIP}
    \end{IF}
    \stm{k:= k+1}
  \end{WHILE}
\end{stuki}

%%% MAXIMUMKERESÉS TÉTEL
\subsection{Maximumkeresés}
Legyen $\mathcal{H}$ egy tetszőleges rendezett halmaz és $f\colon \mathbb{Z}\rightarrow\mathcal{H}$ egy adott függvény. Keressük meg az $f$ függvény maximumát és egy olyan helyét, ahol ezt a maximumértéket felveszi egy adott $[m..n]\subset \mathbb{Z}$ intervallumban.\\

%%% specifikáció
$A = (m\lv\mathbb{Z}, n\lv\mathbb{Z}, i\lv\mathbb{Z}, max\lv\mathcal{H})$\\[2pt]
$B = (m'\lv\mathbb{Z}, n'\lv\mathbb{Z})$\\[2pt]
$Q = (m = m' \land n = n' \land m \leqslant n)$\\[2pt]
$R = (Q \land i \in [m..n] \land max = f(i) \land \forall j \in [m..n]\colon f(j) \leqslant max)$
%%% struktogram
\medskip
\begin{stuki}
  \stm{i,k,max:= m,m,f(m)}
  \begin{WHILE}{3}{\stm{k \neq n}}
  	\begin{CASE}{1}{2}
  		\WHEN{\stm{f(k+1) \geqslant max}}
  		\stm{i,max:= k+1,f(k+1)}
  		\WHEN{\stm{f(k+1) \leqslant max}}
  		\stm{SKIP}      
  	\end{CASE}
    \stm{k:= k+1}
  \end{WHILE}
\end{stuki}

%%% FELTÉTELES MAXIMUMKERESÉS TÉTEL
\subsection{Feltételes maximumkeresés}
Legyen $\mathcal{H}$ egy tetszőleges rendezett halmaz, $f\colon \mathbb{Z}\rightarrow\mathcal{H}$ és $\beta\colon \mathbb{Z}\rightarrow\mathbb{L}$ adott függvények. Határozzuk meg a $\lceil\beta\rceil \cap [m..n]$ halmaz felett az $f$ függvény maximumát és a halmaz egy olyan elemét, amelyen $f$ a maximumértéket felveszi.\\

%%% specifikáció
$A = (m\lv\mathbb{Z}, n\lv\mathbb{Z}, i\lv\mathbb{Z}, max\lv\mathcal{H}, l\lv\mathbb{L})$\\[2pt]
$B = (m'\lv\mathbb{Z}, n'\lv\mathbb{Z})$\\[2pt]
$Q = (m = m' \land n = n' \land m \leqslant n+1)$\\[2pt]
$\begin{aligned}
R& = (Q \land l = (\exists j\in [m..n]\colon \beta(j)) \land l \rightarrow (i \in [m..n] \land \beta(i) \land max = f(i) \land\\
 &\quad\quad\forall j \in [m..n]\colon \beta(j) \rightarrow f(j) \leqslant max))
\end{aligned}$
%$R =  \land $
%%% struktogram
\medskip
\begin{stuki}[17cm]
  \stm{k,l:= m-1,hamis}
  \begin{WHILE}{4}{\stm{k \neq n}}
    \begin{CASE}{2}{7}
	  \WHEN[1]{\stm{\neg\beta(k+1)}}
	    \stm{SKIP}
	  \WHEN[2]{\stm{\beta(k+1) \land \neg l}}
	    \stm{l,i,max:= \\igaz,k+1,f(k+1)}			
	  \WHEN[4]{\stm{\beta(k+1) \land l}}
  	    \begin{CASE}{1}{2}
          \WHEN{\stm{f(k+1) \geqslant max}}
	        \stm{i,max:= k+1,f(k+1)}
	      \WHEN{\stm{f(k+1) \leqslant max}}
	        \stm{SKIP}      
        \end{CASE}			
      \end{CASE}
	  \stm{k:= k+1}
	\end{WHILE}
\end{stuki}

%%% LINEÁRIS KERESÉS TÉTEL - 1. változat
\subsection{Lineáris keresés (1. változat)}
Legyen $\beta\colon \mathbb{Z}\rightarrow\mathbb{L}$ egy adott függvény. Keressük meg azt a legkisebb $\beta$ tulajdonságú egész számot, amely nem kisebb mint az adott $m\in \mathbb{Z}$ szám, feltéve hogy van ilyen.\\

%%% specifikáció
$A = (m\lv\mathbb{Z}, i\lv\mathbb{Z})$\\[2pt]
$B = (m'\lv\mathbb{Z})$\\[2pt]
$Q = (m = m' \land \exists j\in \mathbb{Z}\colon j\geqslant m \land \beta(j))$\\[2pt]
$R = (Q \land i\geqslant m \land \beta(i) \land \forall j \in [m..i-1]\colon \neg\beta(j)))$
%%% struktogram
\medskip
\begin{stuki}[4cm]
  \stm{i:= m}
  \begin{WHILE}{1}{\stm{\neg\beta(i)}}
      \stm{i:= i + 1}
  \end{WHILE}
\end{stuki}
\bigskip

%%% LINEÁRIS KERESÉS TÉTEL - 2.8 változat
\subsection{Lineáris keresés (2.8 változat)}
Legyen $\beta\colon \mathbb{Z}\rightarrow\mathbb{L}$ egy adott függvény. Keressük meg az első olyan helyet az adott $[m..n]$ intervallumban, ahol $\beta$ igaz értéket vesz fel, ha egyáltalán van ilyen hely.\\

%%% specifikáció
$A = (m\lv\mathbb{Z}, n\lv\mathbb{Z}, i\lv\mathbb{Z}, l\lv\mathbb{L})$\\[2pt]
$B = (m'\lv\mathbb{Z}, n'\lv\mathbb{Z})$\\[2pt]
$Q = (m = m' \land n = n' \land m \leqslant n+1)$\\[2pt]
$R = (Q \land l = (\exists j\in [m..n]\colon \beta(j)) \land l \rightarrow (i \in [m..n] \land \beta(i) \land \forall j \in [m..i-1]\colon \neg\beta(j)))$
%%% struktogram
\medskip
\begin{stuki}
  \stm{i,l:= m-1,hamis}
  \begin{WHILE}{2}{\stm{\neg l \land i \le n}}
      \stm{l:= \beta(i+1)}
      \stm{i:= i + 1}
  \end{WHILE}
\end{stuki}

\end{document}